\scanpage{115}
Soit $f(z)=\prod_1^\infty(1+\lambda_i z)$ une fonction entière
de \uline{genre} $1$, posons $\alpha_i=|\lambda_i|$, $F(z)=\prod_{\unsure{1|i}}^\infty(1+\alpha_i z)$.
Alors \uline{$f$ et $F$ ont $=$ ordre} (qui s'exprime donc
aussi en fonction de la suite des $|\alpha_i|$). En effet,
l'ordre de $f\leq$ ordre de $F$ (évident). D'autre
part, ordre de $F=$ ordre de la suite des $|\alpha_i|$
(voir ci-dessus) $\leq$ ordre de $f$ (th. de Hadamard)
\begin{center}\rule{1.5cm}{0.4pt}\end{center}

On peut parler de l'ordre d'un ensemble \inserted{dénombrable} de
nombres complexes, \deleted{\unsure{indépendamment de}} \inserted{$=$ \unsure{non} compris en suite :}
on prend la borne inf. des $s>0$ tels que
$\sum|z_i|^s<+\infty$. Suite à décroissance rapide $=$
d'ordre zéro.

Ordre d'une suite dans un espace de Banach,
ou de Fréchet. (deux définitions, suivant qu'on
prend sommabilité forte, ou sommabilité absolue)

Ordre d'une partie bornée (\deleted{plus petit} \inserted{borne inférieure} des
$s>0$ tels que la partie soit contenue dans l'envelop-\newline
pe convexe cerclée fermée d'une suite d'ordre $s$)
\margin{116}
